\documentclass[10pt]{article}
\usepackage[margin=20mm]{geometry}
\usepackage{amsmath,amssymb,mathtools,fontspec,unicode-math,xurl,hyperref}
\IfFontExistsTF{DejaVu Serif}{\setmainfont{DejaVu Serif}}{\setmainfont{Latin Modern Roman}}
\setmathfont{latinmodern-math.otf}
\setlength{\parindent}{0pt}\setlength{\parskip}{6pt}\sloppy
\title{Technical checks on special unipotent counting and parity reduction}
\author{AI-assisted technical review}
\date{}
\begin{document}
\maketitle
Abstract: Two demonstrated obstructions concern the proof as written: an extra induction and restriction doubles a multiplicity in the quaternionic orthogonal counting argument, and the orientation normalization makes the required compatible flag impossible in odd orthogonal rank. Explicit admissible data and exact independent computations are supplied. A separate notation defect concerns degenerate Weyl-group representations. These findings support an overall Major assessment of the pinned manuscript’s proof; they do not establish that its final representation counts are false. Scientific adjudication remains pending.\par
\section*{Exact source and scope}
The reviewed source is D. Barbasch, J.-J. Ma, B. Sun and C.-B. Zhu, Special unipotent representations of real classical groups: Counting and reduction, arXiv:2205.05266v4, submitted 2 January 2025, 76 pages, SHA-256 prefix `1394994203fb811e`. The associated DOI is 10.4171/JEMS/1609; the publisher records online publication on 28 March 2025. No equivalence with the publisher’s version of record is assumed. All locations below are both PDF and printed page numbers.\par
All 76 pages were read. Detailed checks concern the hypotheses of Theorem 2.1, the parity and relevance conventions, the passage from coherent continuation to multiplicities, the explicit type-A paintings, and the type-\(D^*\) calculation in Proposition 8.6. The proof dependencies through Corollary 8.7 and Section 9 were followed. The cell theory, Kazhdan–Lusztig–Vogan algorithm and every cited external theorem were not independently re-proved. The descent cases in Section 10 were read and their role traced, but not exhaustively certified in arbitrary rank.\par
The distinction between an unconditional inequality and a conditional equality in Theorem 2.1 is respected: the equality requires the stated restriction on exceptional factors and linearity or the real metaplectic case. The proof through Propositions 4.17, 4.18, Theorem 4.23 and Corollary 4.24 does not justify dropping these conditions. Genuine representations, the metaplectic dual, and the two very-even orbits must likewise retain their separate conventions.\par
\section*{A doubled multiplicity in the core counting proof}
Problem 1 (Major: Proposition 8.6, p. 59).  In the \(D^*\) computation, the definition \[\mathcal C_b=\operatorname{Ind}_{H_{n_b/2}}^{W_{n_b}}\eta,
 \qquad \mathcal C_b\text{ then restricted to }W'_{n_b},\] is incompatible with the subsequent multiplicity calculation. Proposition 8.1 (p. 53) instead uses \(\operatorname{Ind}_{H_{n_b/2}}^{W'_{n_b}}\eta\). The printed choice introduces an extra copy in an admissible example. In particular the last equality of the displayed calculation on p. 59 is false with its printed data.\par
\subsection*{Admissible orbit and paintings}
Take \(G=SO^*(8)\), \(\star=D^*\), and the dual nilpotent orbit with row partition \[\check{\mathcal O}=(3,2,2,1),\quad
 \check{\mathcal O}_b=(2,2),\quad
 \check{\mathcal O}'_b=(2),\quad
 \check{\mathcal O}_g=(3,1).\] The only even part has even multiplicity, so this is a valid orthogonal nilpotent partition. It is not very even. Here \(n_b=n_g=2\), the good rows are odd, and the orbit is \(G\)-relevant by Proposition 2.19. The type-I choice for the bad part is allowed by the orientation convention discussed at the start of Section 9; alternatively the opposite choice gives the same multiplicity discrepancy.\par
Equations (8.6)–(8.7) give the bad bipartition \(((1),(1))\). Definition 8.5 permits only bullets in both diagrams, so there is exactly one bad painting. For the good partition \((3,1)\), (8.9) gives a left diagram consisting of a column of height two and an empty right diagram. Definition 2.24 permits only \(\bullet,s\) on the left and \(\bullet,r\) on the right; equality of the bullet sets forces both left boxes to be \(s\). This is a valid painting and is unique. There are no primitive pairs. Thus the relevant product of painting cardinalities is \(1\).\par
\subsection*{Exact character calculation}
Equation reference: double. Write \(B=W_2\), the eight signed permutation matrices of size two, and \[D=W'_2=H_1=\{1,s,z,zs\},\qquad
 s=\begin{pmatrix}0&1\\1&0\end{pmatrix},\qquad z=-I_2.\] The definition (8.2) gives \(\eta(s)=1\), \(\eta(z)=-1\). Let \(\rho\) be the natural two-dimensional representation of \(B\), and let \(w=\operatorname{diag}(-1,1)\). Then \[\operatorname{Ind}_D^B\eta=\rho,\qquad
 \operatorname{Res}_D^B\rho=\eta\oplus\eta^w,
 \qquad \eta^w(s)=-1,\quad\eta^w(z)=-1.\] These identities follow either by the four-element character table or directly by diagonalizing \(s\) on \(\mathbb C^2\). Consequently the printed definition gives \[
 \left[\rho:\operatorname{Ind}_D^B\mathcal C_b\right]
 =\left[\rho:\operatorname{Ind}_D^B(\eta\oplus\eta^w)\right]
 =[\rho:2\rho]=2.\] The bad painting number is \(1\), not \(2\).\par
For completeness the good factor is nonzero, so the error does not disappear upon multiplication. Let \(\kappa(s)=-1\), \(\kappa(z)=1\). The good representation \(\widetilde\tau_g\) is the inflation of \(\operatorname{sgn}_{S_2}\) to \(B\), with restriction \(\kappa\). Proposition 8.2 yields \[\mathcal C_g=\eta+\operatorname{Ind}_{\langle s\rangle}^{D}\operatorname{sgn},\qquad
 [\kappa:\mathcal C_g]=1.\] Thus the penultimate product on p. 59, with the printed \(\mathcal C_b\), equals \(2\cdot1=2\), while its claimed painting product equals \(1\).\par
The first equality is affected as well. Put \[L=D\times D,\qquad
 W''=\{(a,b)\in B\times B:\epsilon(a)\epsilon(b)=1\}.\] Then \(|L|=16\) and \(|W''|=32\). Corollary 6.15 and Proposition 8.1 give the coherent module \(\operatorname{Ind}_L^{W''}(\eta\otimes\mathcal C_g)\). Its restriction to \(L\) has multiplicity one at \(\eta\otimes\kappa\). Replacing \(\eta\) by the printed \(\operatorname{Res}_D^B\operatorname{Ind}_D^B\eta\) makes this multiplicity two. This is also obtained by the elementary index-two restriction formula, since \(\kappa\) is fixed by the outer twist while \(\eta\) and \(\eta^w\) are distinct.\par
An original program enumerated all eight matrices, evaluated the induced characters by their defining finite sums, and checked the \(32\)-element group \(W''\). It returned respectively \(1\) and \(2\) for the correct and printed multiplicities; no author implementation was used. These exact calculations corroborate, rather than replace, the analytic argument above.\par
\subsection*{Impact and repair boundary}
Proposition 8.6 is the explicit multiplicity-to-painting step. It feeds Corollary 8.7, the counting input to Proposition 9.1 and the reduction Theorem 2.20; its good-parity specialization is also part of the route to Theorem 2.27. Hence the exhibited equality is inside a core derivation. The example concerns a mixed-parity \(D^*\) case and does not refute Theorem 2.27’s good-parity numerical claim.\par
Replacing \(\mathcal C_b\) by \(\operatorname{Ind}_{H_{n_b/2}}^{W'_{n_b}}\eta\) removes this doubling and restores the checked equality. This is a concrete local repair, not a certification of all ranks or all orbit labels. The degenerate representation conventions discussed in Section 4 below also require consistent correction. The Major assessment rests on the demonstrated invalid core calculation; no claim of an incorrect final packet cardinality is made.\par
\section*{Incompatible orientation and flag normalization}
Problem 2 (Major: (6.4), the condition following (6.7), and Remark 6.2).  For \(\star=D\) or \(D^*\) and odd rank \(n\), the unit orientation condition on p. 40 is incompatible with the compatible-basis condition on p. 41. Consequently Remark 6.2(a)–(b), asserting existence of a compatible isotropic subspace and flag, is false as written.\par
Equation reference: sign. The induced symmetric bilinear form on \(\bigwedge^{2n}V\) is explicitly defined by the determinant on p. 40. Equation (6.4) requires \(\langle\omega_V,\omega_V\rangle=1\). On p. 41 compatibility requires \[\omega_V=u_1\wedge\cdots\wedge u_{2n},\qquad
 \langle u_i,u_j\rangle=\delta_{i+j,\,2n+1}.\] The latter Gram matrix is the reversal matrix \(J_{2n}\), for which \[
 \det J_{2n}=(-1)^{(2n)(2n-1)/2}=(-1)^n.\] For odd \(n\) this forces the same squared norm to be \(-1\).\par
Here is a complete admissible instance avoiding low-rank exceptions. Let \(\star=D\), \(n=5\), \(V=\mathbb C^{10}\), with \(\langle x,y\rangle=x^tJ_{10}y\), \(j(x)=\overline{x}\), and \[\omega_V=i\,e_1\wedge\cdots\wedge e_{10}.\] The form is nondegenerate and symmetric, \(j^2=1\), and (6.3) holds because \(J_{10}\) is real. Moreover \(\langle\omega_V,\omega_V\rangle=i^2\det J_{10}=1\). The \(D^*\)-only condition (6.5) and the dimension-zero condition (6.6) do not apply. Thus this is an allowed \([D]\)-space, with \(G(V)=SO(5,5)\). If a nilpotent orbit is needed in the subsequent setup, choose the zero dual orbit \((1^{10})\); it has good parity, \(n_b=0\), and satisfies relevance.\par
Every basis with the prescribed Gram matrix has wedge norm \(-1\), so none can have wedge equal to this \(\omega_V\). Equivalently, a matrix \(U\) satisfying \(U^tJ_{10}U=J_{10}\) has \((\det U)^2=1\), whereas its wedge would need \(\det U=i\). There is no compatible maximal isotropic subspace in the stated sense, and hence no flag of the required kind. Symbolic determinants for ranks \(1\) through \(7\) independently confirm [sign] (\#sign).\par
The proof uses that flag to obtain (6.8), then the coordinates, positive roots, parity splitting and parameter matching in Sections 6.2–6.5. This is an obstruction to the stated construction for an infinite family of allowed groups, rather than an example violating a hypothesis. It does not establish failure of the intrinsic representation-theoretic results. A possible local convention is \(\omega_V=c_n\,u_1\wedge\cdots\wedge u_{2n}\) with \(c_n^2=(-1)^n\), for example \(c_n=i^n\). One must then reconcile this choice with (6.5), direct products, dual orientations and the type-I/II labels; the present review does not certify that global repair.\par
\section*{Degenerate-label notation}
Problem 3 (Minor: Section 8.3, p. 55, and the recall on p. 59).  The printed definition of \((\tau_L,\tau_R)_I\) for \(\tau_L=\tau_R\) uses \(\operatorname{Ind}_{S_t}^{W'_t}\tau_L\). But \(|\tau_L|=t/2\), so \(\tau_L\) labels an \(S_{t/2}\)-representation, not an \(S_t\)-representation under the paper’s preceding convention.\par
For \(t=2\) and \(\tau_L=\tau_R=(1)\) the stated input is not an \(S_2\)-module as defined. Even if one supplies the trivial \(S_2\)-module by hand, its ordinary induction to \(W'_2\) has dimension two and is reducible; the two constituents of the restricted natural \(W_2\) representation are one-dimensional. Thus ordinary induction cannot be the claimed irreducible constituent. A consistent choice of degenerate constituents, and its relation to the relevant branching formula (8.16), must be specified.\par
The p. 59 recall \(\widetilde\tau_b=(\mathcal O'_b,\mathcal O'_b)\in\operatorname{Irr}(W_{n_b})\) has a related size mismatch: \(|\mathcal O'_b|=n_b\), so that bipartition has total size \(2n_b\). In the example of Section 2 it would label a representation of \(W_4\), not \(W_2\). The earlier \((\tau_{L,b},\tau_{R,b})\) has the required size. These are supported local notation defects. No additional numerical failure independent of Problem 1 is inferred from an unspecified interpretation, which is why this finding is classified separately as Minor.\par
\section*{Other checks and limits}
The good-parity test can be checked directly on an \(\mathfrak{sl}_2\) block of dimension \(r\): the eigenvalues of its neutral element are \(r-1,r-3,\ldots,1-r\), so \(\exp(\pi i h)\) acts as \((-1)^{r-1}\). Comparing this scalar to that of the principal block recovers the parities in Definition 2.6; passage to the quotient by \(\{\pm I\}\) explains the different all-rows-same-parity statement for \(\widetilde A\) in Lemma 2.7.\par
An independent exhaustive enumeration applied the actual prefix, row and column rules of Definition 2.8 to all \(67\) partitions of total size at most eight, including the empty partition. For \(A_{\mathbb R}\) it recovered \(\prod_r(1+m_r)\), with \(m_r\) the row multiplicity in the dual partition. For \(A_{\mathbb H}\) it recovered one painting exactly when every row length is even. The two shapes in Example 2.28 each gave four paintings, consistent with the stated count eight after the primitive-pair factor. These bounded checks do not prove all of the descent recursions or the irreducibility of induced representations.\par
The exact source pages containing the challenged formulas were actually viewed, including pp. 40–42, 55 and 59, so none of the findings rests only on a lost prime, overbar or diagram in extracted text. The checked proof dependencies are distinct from unverified external inputs. The overall assessment is Major, with the scope and repair limits above; this initial review is not independent acceptance.\par
\section*{Related materials}
The author publication page links a Python repository to the separate construction and unitarity paper, not to this counting-and-reduction v4. Its public README was inspected; no version-matched release for the present manuscript was established, and no author code was executed. Bounded title/identifier searches found no dedicated correction or supplement. This is a search limitation, not evidence of absence or a mathematical defect. The pinned publisher and arXiv metadata remain the identity evidence.\par
9 D. Barbasch, J.-J. Ma, B. Sun and C.-B. Zhu, Special unipotent representations of real classical groups: Counting and reduction, arXiv:2205.05266v4 (2 January 2025). Associated DOI: 10.4171/JEMS/1609 (https://doi.org/10.4171/JEMS/1609). J.-J. Ma, Publications/preprints, https://www.majiajun.org/publication/, consulted 27 October 2026. J.-J. Ma, unipotentrepn, public repository README, https://github.com/jiajunma/unipotentrepn, consulted 27 October 2026. The repository is linked by the author to the companion construction paper.\par
\end{document}
